\chapter{Fano--Feshbach 共振与 Landau--Zener 跃迁}\label{ch:resonance-lz}

共振问题有两类常被分开讲：一类是离散态嵌入连续谱形成 Fano/Feshbach resonance，另一类是两个能级随参数变化形成 avoided crossing 并发生 Landau--Zener 跃迁。它们的共同点是：完整 Hilbert 空间中存在一个较小的“关注子空间”，其余自由度通过能量依赖的自能或时间依赖的有效耦合反馈回来。本章先建立一般投影理论，再分别进入定态共振与含时 crossing 两个可解模型。

\section{Feshbach 投影：先消去不关注的子空间}

设正交投影算符

\[
P^2=P,
\qquad
Q^2=Q,
\qquad
PQ=0,
\qquad
P+Q=\mathbb I.
\]

更一般地，对复谱参数 $z$，完整 resolvent $G(z)=(z-H)^{-1}$ 的 $P$ 块满足精确 Schur-complement 恒等式
\begin{equation}
\boxed{
PG(z)P
=\left[z-H_{PP}-H_{PQ}(z-H_{QQ})^{-1}H_{QP}\right]^{-1}_{P}.
}
\label{eq:feshbach-schur-resolvent}
\end{equation}
这里右端逆只在 $P$ 子空间中取。定态有效 Hamiltonian、共振极点、能级位移和寿命都只是这一 resolvent 恒等式在不同解析区域与边界条件下的读法。


\begin{derivation}[Schur complement 如何给出 Feshbach resolvent]
设 Hilbert 空间分解 \(\mathcal H=\mathcal H_P\oplus\mathcal H_Q\)，正交投影 \(P+Q=I\)。对 \(z\in\rho(H)\)，resolvent \(G(z)=(z-H)^{-1}\) 满足
\[
 (z-H)|\Psi\rangle=|\eta_P\rangle,
 \qquad Q|\eta_P\rangle=0.
\]
把 \(|\Psi\rangle\) 分块为 \(|\Psi\rangle=|\psi_P\rangle+|\psi_Q\rangle\)，其中 \(|\psi_P\rangle=P|\Psi\rangle\)、\(|\psi_Q\rangle=Q|\Psi\rangle\)。将方程分别左乘 \(P\) 和 \(Q\)，利用 \(PHP=H_{PP}\)、\(PHQ=H_{PQ}\) 等分块记号，得到耦合方程组
\begin{align*}
 (z-H_{PP})|\psi_P\rangle-H_{PQ}|\psi_Q\rangle
 &=|\eta_P\rangle,\\
 -H_{QP}|\psi_P\rangle+(z-H_{QQ})|\psi_Q\rangle
 &=0.
\end{align*}
第一式是 \(P\) 空间分量，第二式是 \(Q\) 空间分量，两式通过非对角块 \(H_{PQ}\)、\(H_{QP}\) 耦合。

只要 \(z\) 位于 \(H_{QQ}\) 的 resolvent set，\((z-H_{QQ})\) 可逆，第二式精确解为
\[
 |\psi_Q\rangle
 =(z-H_{QQ})^{-1}H_{QP}|\psi_P\rangle.
\]
这一步把 \(Q\) 空间分量完全消去：给定 \(|\psi_P\rangle\)，\(|\psi_Q\rangle\) 唯一确定。物理上，\(Q\) 空间被"积分掉"，其效应完全编码在有效 Hamiltonian 中。

将 \(|\psi_Q\rangle\) 代回第一式，
\[
 (z-H_{PP})|\psi_P\rangle
 -H_{PQ}(z-H_{QQ})^{-1}H_{QP}|\psi_P\rangle
 =|\eta_P\rangle,
\]
整理为
\[
 \left[z-H_{PP}
 -H_{PQ}(z-H_{QQ})^{-1}H_{QP}\right]|\psi_P\rangle
 =|\eta_P\rangle.
\]
方括号中正是 Feshbach 有效 Hamiltonian
\[
 H_{\rm eff}(z)=H_{PP}+H_{PQ}(z-H_{QQ})^{-1}H_{QP},
\]
因此 \(|\psi_P\rangle=(z-H_{\rm eff}(z))^{-1}|\eta_P\rangle\)。另一方面，由 \(|\eta_P\rangle\in\mathcal H_P\) 和 \(G(z)=(z-H)^{-1}\)，
\[
 |\psi_P\rangle=P|\Psi\rangle=PG(z)|\eta_P\rangle=PG(z)P|\eta_P\rangle.
\]
两个表达式对任意 \(|\eta_P\rangle\in\mathcal H_P\) 成立，故算子恒等
\[
 PG(z)P=(z-H_{\rm eff}(z))^{-1}
 =\left[z-H_{PP}-H_{PQ}(z-H_{QQ})^{-1}H_{QP}\right]^{-1},
\]
即\eqnrefs{eq:feshbach-schur-resolvent}。

有效 Hamiltonian 的结构有三层物理含义。\(H_{PP}\) 是 \(P\) 子空间中的裸 Hamiltonian；\((z-H_{QQ})^{-1}\) 是 \(Q\) 空间的 restricted resolvent，描述中间态在 \(Q\) 空间中传播；\(H_{PQ}\) 和 \(H_{QP}\) 分别是 \(P\to Q\) 和 \(Q\to P\) 的耦合矩阵元。整个修正项 \(H_{PQ}(z-H_{QQ})^{-1}H_{QP}\) 是"从 \(P\) 跳到 \(Q\)、在 \(Q\) 中传播、再跳回 \(P\)"的二阶虚过程。

当 \(z\) 逼近 \(Q\) 空间连续谱时，\((z-H_{QQ})^{-1}\) 不再有界。此时需指定边界条件：出射波对应 \(z\to E+\ii0^+\)，入射波对应 \(z\to E-\ii0^+\)。由 Sokhotski--Plemelj 公式，
\[
 (E\pm\ii0^+-H_{QQ})^{-1}
 =\operatorname{p.v.}\frac1{E-H_{QQ}}
 \mp\ii\pi\delta(E-H_{QQ}),
\]
虚部产生宽度 \(\Gamma=-2\operatorname{Im}\langle\psi_P|H_{PQ}\delta(E-H_{QQ})H_{QP}|\psi_P\rangle\)，即 Fermi 黄金法则中的衰变率。有效 Hamiltonian 的非 Hermitian 部分正由这一步产生，把实能级移动到复平面 \(E\to E-\ii\Gamma/2\)，给出指数衰变 \(e^{-\Gamma t/2}\)。
\end{derivation}

定态 Schr\"odinger 方程

\[
(E-H)|\Psi\rangle=0
\]

分块为

\begin{align*}
(E-H_{PP})P|\Psi\rangle
&=H_{PQ}Q|\Psi\rangle,
\\
(E-H_{QQ})Q|\Psi\rangle
&=H_{QP}P|\Psi\rangle.
\end{align*}

散射问题要求 outgoing boundary condition，因此 resolvent 应取

\[
G_Q^{(+)}(E)
\equiv(E+\ii0^+-H_{QQ})^{-1}.
\]

第二式形式精确地解为

\[
Q|\Psi\rangle
=G_Q^{(+)}(E)H_{QP}P|\Psi\rangle.
\]

代回 \(P\) 空间：

\begin{equation}
\boxed{
H_{\rm eff}^{(+)}(E)
=H_{PP}
+H_{PQ}G_Q^{(+)}(E)H_{QP}.
}
\label{eq:feshbach-effective}
\end{equation}

第二项

\[
\Sigma(E)
\equiv H_{PQ}G_Q^{(+)}(E)H_{QP}
\]

就是能量依赖自能。它的实部移动共振位置，虚部产生有限寿命。只要 \(Q\) 空间包含连续谱，\(H_{\rm eff}^{(+)}\) 一般就是非 Hermitian 的；这不是基本 Hamiltonian 失去 Hermiticity，而是 outgoing boundary condition 和连续谱消元后的有效描述。

\section{定态共振：自能、Fano 干涉与阈值散射}
接下来考察单离散态耦合连续谱。

取

\[
P=|\phi\rangle\langle\phi|,
\qquad
H_0|\phi\rangle=E_\phi|\phi\rangle,
\]

并令连续态满足能量归一化

\[
\langle E|E'\rangle=\delta(E-E').
\]

定义耦合

\[
V_E\equiv\langle E|V|\phi\rangle.
\]

则标量自能为

\[
\Sigma(E)
=\int\dd E'\,
\frac{|V_{E'}|^2}{E-E'+\ii0^+}.
\]

Sokhotski--Plemelj 公式

\[
\frac1{x+\ii0^+}
=\mathcal P\frac1x-\ii\pi\delta(x)
\]

给出

\begin{equation}
\boxed{
\Sigma(E)
=\Delta(E)-\frac{\ii}{2}\Gamma(E),
}
\label{eq:self-energy-split}
\end{equation}

其中

\begin{equation}
\boxed{
\Delta(E)
=\mathcal P\int\dd E'\,
\frac{|V_{E'}|^2}{E-E'},
\qquad
\Gamma(E)=2\pi|V_E|^2.
}
\label{eq:shift-width}
\end{equation}

\(\Gamma\) 的具体前因子依赖连续态归一化；如果改用动量归一化，态密度 Jacobian 必须显式出现。

共振极点由

\[
E-E_\phi-\Sigma(E)=0
\]

决定。若 \(\Delta(E),\Gamma(E)\) 在共振附近缓慢变化，定义

\[
E_R=E_\phi+\Delta(E_R),
\]

就得到近似极点

\[
E_{\rm pole}\simeq E_R-\ii\Gamma(E_R)/2.
\]

对应时间振幅 \(\exp(-\ii E_{\rm pole}t/\hbar)\) 的概率按 \(\exp(-\Gamma t/\hbar)\) 衰减，所以寿命

\[
\tau\simeq\frac{\hbar}{\Gamma}.
\]

若自能的能量依赖不能忽略，极点附近传播子不是简单的单位 residue。设 $z_p-E_\phi-\Sigma(z_p)=0$，则
\begin{equation}
G_P(z)\simeq\frac{Z_p}{z-z_p},
\qquad
Z_p=\frac{1}{1-\Sigma'(z_p)}.
\label{eq:resonance-pole-residue}
\end{equation}
$Z_p$ 衡量裸离散态在 dressed resonance 中的极点权重；宽共振或阈值附近 $\Sigma'(z)$ 较大时，把 $Z_p$ 近似为 1 会丢失重要谱权重信息。

接下来考察Fano 散射态的显式构造。

仅知道复极点还不能解释不对称 Fano profile，因为不对称来自“背景连续路径”和“经离散态再回到连续谱的路径”相干叠加。设固定总能量为 \(E\) 的 exact eigenstate 展开为

\begin{equation}
|\Psi_E\rangle
=a_E|\phi\rangle
+\int\dd E'\,b_E(E')|E'\rangle.
\label{eq:fano-state-expansion}
\end{equation}

Schr\"odinger 方程给出

\begin{align*}
(E-E_\phi)a_E
&=\int\dd E'\,V_{E'}^*b_E(E'),
\\
(E-E')b_E(E')
&=V_{E'}a_E.
\end{align*}

第二式在分布意义下的一般解不是普通函数 \(V_{E'}a_E/(E-E')\)，因为在 \(E'=E\) 处还允许 delta 函数齐次解。写成

\begin{equation}
\boxed{
b_E(E')
=a_EV_{E'}
\left[
\mathcal P\frac1{E-E'}
+Z_E\delta(E-E')
\right].
}
\label{eq:fano-distribution-solution}
\end{equation}

代回第一式：

\[
E-E_\phi
=\Delta(E)+Z_E|V_E|^2.
\]

引入

\[
\epsilon
\equiv\frac{E-E_R}{\Gamma/2},
\qquad
\Gamma=2\pi|V_E|^2,
\]

得到

\[
\boxed{Z_E=\pi\epsilon.}
\]

为了按能量 delta 函数归一化，可定义散射相移 \(\delta_E\) 使

\[
Z_E=-\pi\cot\delta_E.
\]

于是

\begin{equation}
\boxed{
\tan\delta_E=-\frac1\epsilon.
}
\label{eq:fano-phase-shift}
\end{equation}

同时选择

\[
a_E=\frac{\sin\delta_E}{\pi V_E}
\]

（若 \(V_E\) 为复数，相位可吸收到连续态 convention 中）。定义 dressed discrete state

\[
|\Phi_E\rangle
=|\phi\rangle
+\mathcal P\int\dd E'\,
\frac{V_{E'}}{E-E'}|E'\rangle,
\]

则 exact scattering state 变为

\begin{equation}
\boxed{
|\Psi_E\rangle
=\frac{\sin\delta_E}{\pi V_E}|\Phi_E\rangle
-\cos\delta_E|E\rangle.
}
\label{eq:fano-normalized-state}
\end{equation}

这一步把“连续背景”和“共振 dressing”两条振幅路径显式分开了。

\begin{derivation}[Fano 线型]
\emph{第一步：跃迁振幅的两路径分解。}设外部跃迁算符为 \(T\)，初态为 \(|i\rangle\)。由 \eqnrefs{eq:fano-normalized-state}，散射态 $|\Psi_E\rangle$ 已分解为连续背景 $|E\rangle$ 和 dressed 离散态 $|\Phi_E\rangle$ 两条路径，跃迁振幅为
\[
\mathcal M(E)
=\langle\Psi_E|T|i\rangle
=-\cos\delta_E\,\mathcal M_{\rm bg}
+\frac{\sin\delta_E}{\pi V_E^*}
\langle\Phi_E|T|i\rangle,
\]
其中 \(\mathcal M_{\rm bg}=\langle E|T|i\rangle\) 是纯背景路径的振幅。两条路径的相干叠加是不对称线型的根源。

\emph{第二步：定义 Fano $q$ 参数。}定义
\[
q(E)
\equiv
\frac{\langle\Phi_E|T|i\rangle}
{\pi V_E^*\mathcal M_{\rm bg}}.
\]
$q$ 是离散路径与背景路径的约化振幅比，包含两条路径的相对相位和强度。于是
\[
\frac{\mathcal M}{\mathcal M_{\rm bg}}
=-\cos\delta_E+q\sin\delta_E.
\]

\emph{第三步：用相移表达。}由 \(\tan\delta_E=-1/\epsilon\)，可取
\[
\cos\delta_E=\frac{\epsilon}{\sqrt{1+\epsilon^2}},
\qquad
\sin\delta_E=-\frac1{\sqrt{1+\epsilon^2}},
\]
代入得
\[
\frac{\mathcal M}{\mathcal M_{\rm bg}}
=\frac{q-\epsilon}{\sqrt{1+\epsilon^2}},
\]
整体负号不影响截面，因此
\[
\boxed{
\frac{|\mathcal M(E)|^2}{|\mathcal M_{\rm bg}(E)|^2}
=\frac{(q+\epsilon)^2}{1+\epsilon^2}.
}
\]
不对称性由两条复振幅的相对大小和相位产生，而不是 Lorentzian 人为"加一个偏斜参数"。

\emph{第四步：极限线型。}$|q|\to\infty$ 时，$(q+\epsilon)^2/(1+\epsilon^2)\approx q^2/(1+\epsilon^2)$，退化为以 $\epsilon=0$ 为中心的 Lorentz 峰（离散路径主导，背景可忽略）。$q=0$ 时，$\epsilon^2/(1+\epsilon^2)$ 在 $\epsilon=0$ 处为零、$|\epsilon|\to\infty$ 处趋于 1，共振中心成为反共振零点（完全相消干涉）。$q=\pm1$ 时，$(1+\epsilon)^2/(1+\epsilon^2)$ 在 $\epsilon=-1$ 处为零、$\epsilon=0$ 处为 $1/2$，不对称最显著。

\emph{举例：氦原子自电离共振。}氦的 $2s2p\,^1P$ 自电离态耦合到 $1s\epsilon p\,^1P$ 连续态，$I_p\approx60\,$eV 处的共振 $q\approx-2.8$、$\Gamma\approx0.14\,$eV。光吸收截面在共振低能侧（$\epsilon\approx-2.8$）出现尖锐零点，高能侧呈 Lorentz 尾。这一不对称 profile 直接来自离散态 $2s2p$ 与连续态 $1s\epsilon p$ 两条光吸收路径的 Fano 干涉，$q$ 的具体值由两条路径的偶极矩阵元比决定。

\emph{更一般结论：多通道 Fano。}若多个离散态 $\{|\phi_n\rangle\}$ 同时耦合到同一连续谱，有效 Hamiltonian 是 $N\times N$ 矩阵 $H_{mn}=E_n\delta_{mn}+\Delta_{mn}(E)-\ii\pi V_{nE}V_{mE}^*/2$，Fano profile 推广为多共振叠加，$q$ 参数变为矩阵值。若连续谱有多个通道（多通道散射），$q$ 一般为复数，其虚部来自通道间吸收或非弹性耦合。这些推广都建立在同一 Schur 补消元结构 \eqnrefs{eq:feshbach-schur-resolvent} 之上。

\emph{物理意义。}Fano 线型的物理本质是"离散共振态与连续背景的量子干涉"：跃迁振幅分解为两条路径——经离散态的共振路径和经连续态的背景路径，截面 \(|\mathcal M|^2\) 中两条路径的交叉项产生不对称线型。Fano 参数 \(q\) 是两条路径的约化振幅比，编码了相对相位和强度：\(|q|\to\infty\) 时离散路径主导，退化为对称 Lorentz 峰（标准共振）；\(q=0\) 时完全相消干涉，共振中心成为反共振零点（"暗态"）；\(q\sim\pm1\) 时不对称最显著。与 Breit--Wigner 共振的根本区别在于连续背景的存在：Breit--Wigner 假设离散态与连续态解耦，截面是纯 Lorentzian；Fano 保留耦合，干涉零点 \(\epsilon=-q\) 出现在共振附近，其位置直接给出 \(q\)。在实验上，Fano 线型普遍出现于自电离共振（氦 \(2s2p\)）、Rydberg 态自电离、量子点输运和光子晶体缺陷模等任何"离散态嵌入连续谱"的体系。
\end{derivation}

因此 Fano profile 为

\begin{equation}
\boxed{
\sigma(E)
=\sigma_{\rm bg}(E)
\frac{(q+\epsilon)^2}{1+\epsilon^2}.
}
\label{eq:fano-profile}
\end{equation}

在单通道、时间反演对称且无吸收的常见 convention 中，可以通过相位选择把 $q$ 取为实数，得到\eqnrefs{eq:fano-profile}。若存在不可消去的多通道相位、吸收或时间反演破缺，$q$ 一般可以是复数，此时线型应保留 $|q+\epsilon|^2/(1+\epsilon^2)$，不能强行用实 $q$ 描述所有不对称性。

不同 \(q\) 的线型见\figref{fig:fano-lineshape}。

\begin{figure}[tbp]
\centering
\includegraphics[width=.83\textwidth]{fano_lineshape.pdf}
\caption{不同 Fano 参数 \(q\) 对归一化线型的影响。\(|q|\gg1\) 时共振趋向 Lorentz 峰，\(q=0\) 时共振中心成为反共振零点。}
\label{fig:fano-lineshape}
\end{figure}

接下来考察Feshbach 共振的低能散射极限。

冷原子 Feshbach resonance 使用的是同一投影结构，只是空间含义对调：\(P\) 是开放散射通道，\(Q\) 包含闭通道束缚态。若闭通道中只有一个束缚态 \(|\phi_c\rangle\) 靠近开放通道阈值，则开通道的有效相互作用包含

\[
V_{\rm eff}(E,B)
\supset
\frac{W|\phi_c\rangle\langle\phi_c|W^\dagger}
{E-E_c(B)-\Sigma_c(E)}.
\]

在超低能 s 波区，散射振幅满足有效程展开

\[
k\cot\delta_0(k)
=-\frac1a+\frac12r_ek^2+\cdots.
\]

若闭通道束缚态相对于阈值的 detuning 在窄磁场区间近似线性，

\[
E_c(B)-E_{\rm th}
\simeq\delta\mu(B-B_c),
\]

并把实自能产生的阈值位移吸收到新的共振位置 \(B_0\)，则零能 \(T\)-matrix 中的共振项必然产生一个 \((B-B_0)^{-1}\) 极点。因此散射长度可以参数化为

\begin{equation}
\boxed{
a(B)
=a_{\rm bg}
\left(1-\frac{\Delta B}{B-B_0}\right).
}
\label{eq:feshbach-scattering-length}
\end{equation}

\(a_{\rm bg}\) 吸收所有缓慢变化的开放通道背景，\(\Delta B\) 把耦合强度和磁矩差压缩成实验可测的磁场宽度。\eqnrefs{eq:feshbach-scattering-length} 是“单闭通道束缚态 + 低能 s 波 + 背景缓慢变化”的阈值近似，而不是一般 Feshbach 理论的定义。


\begin{derivation}[从闭通道极点到散射长度的磁场极点，即\eqnrefs{eq:feshbach-scattering-length}]
\emph{物理设置。}Feshbach 共振中，开通道与闭通道通过磁场耦合。闭通道束缚态能量 $E_c(B)$ 随磁场线性移动：$E_c(B)\approx\delta\mu(B-B_c)$，其中 $\delta\mu$ 是开闭通道磁矩差，$B_c$ 是裸闭通道态与阈值交叉的位置。$\delta\mu$ 的典型量级是核磁矩 $\mu_N$，远小于电子磁矩 $\mu_B$，故磁场调谐精度高。

\emph{第一步：消去闭通道。}在零碰撞能附近，把开通道背景和耦合矩阵元视为对 $B$ 缓慢变化。单闭通道态消元（Feshbach 投影 $P$-$Q$ 分割）后，零能散射振幅的共振部分具有结构
\begin{equation*}
 f_{\rm res}(0,B)
 =\frac{C_f}
 {-\delta\mu(B-B_c)-\Delta_c(0)},
\end{equation*}
其中 $C_f$ 汇总开--闭通道耦合和归一化常数，$\Delta_c(0)$ 是实自能位移（虚部在零能为零，因为开通道 s 波阈值处没有连续态可衰减）。这一步直接应用了 \eqnrefs{eq:feshbach-effective} 的有效 Hamiltonian，闭通道的 $(z-H_{QQ})^{-1}$ 在单束缚态近似下化为 $1/(E-E_c(B)-\Delta_c(E))$。

\emph{第二步：定义物理共振位置。}极点发生在分母为零处。定义物理共振位置
\begin{equation*}
 B_0
 =B_c-\frac{\Delta_c(0)}{\delta\mu},
\end{equation*}
即实自能把裸交叉 $B_c$ 位移到 $B_0$。$B_0$ 与 $B_c$ 的差 $\Delta_c(0)/\delta\mu$ 来自闭通道实自能，包含开通道连续态的虚跃迁贡献。于是散射振幅分解为背景加共振极点：
\begin{equation*}
 f(0,B)
 =f_{\rm bg}(0)-\frac{C_f/\delta\mu}{B-B_0}.
\end{equation*}
背景 $f_{\rm bg}$ 吸收所有非共振的开通道散射信息，共振极点是唯一的磁场依赖结构。在 $B=B_0$ 附近，共振项主导，散射振幅发散；远离 $B_0$ 时，背景项主导，$f\approx f_{\rm bg}$。

\emph{第三步：散射长度参数化。}三维 s 波阈值定义为 $f_0(k\to0)=-a$，于是
\begin{equation}
 a(B)=a_{\rm bg}-\frac{C_a}{B-B_0},
 \qquad
 C_a\equiv-\frac{C_f}{\delta\mu}.
 \label{eq:feshbach-scattering-length-additive}
\end{equation}
这是加法形式\eqnrefs{eq:feshbach-scattering-length-additive}。$C_a$ 的符号由开闭通道耦合的相位约定决定，实验上通过测量 $a(B)$ 的发散方向确定。

\emph{加法与乘法形式的适用范围。}加法形式 \eqnrefs{eq:feshbach-scattering-length-additive} 在 $a_{\rm bg}=0$ 时仍有效，而乘法形式 \eqnrefs{eq:feshbach-scattering-length} 此时退化。反之，当 $|a_{\rm bg}|\gg|C_a/(B-B_0)|$ 时乘法形式更自然，因为它显式分离了背景与共振。两种形式的等价性来自 $a_{\rm bg}\neq0$ 的假设；实验上若发现 $a(B)$ 在远离共振处接近零，应改用加法形式拟合。

\emph{第四步：乘法形式。}只要 $a_{\rm bg}\neq0$，定义 $\Delta B=C_a/a_{\rm bg}$，则
\begin{equation*}
 a(B)=a_{\rm bg}\left(1-\frac{\Delta B}{B-B_0}\right),
\end{equation*}
即\eqnrefs{eq:feshbach-scattering-length}。因此乘法形式不是比加法形式更基本；当背景散射长度接近零时，\eqnrefs{eq:feshbach-scattering-length-additive}反而是更稳健的参数化。$a(B)$ 在 $B=B_0$ 处发散（共振极点），在 $B=B_0+\Delta B$ 处过零（zero crossing），这两个特征是实验标定 Feshbach 共振的主要依据。

\emph{共振宽度的物理含义。}$\Delta B$ 不是共振"强度"而是"宽度"：$|\Delta B|$ 大表示共振影响范围广，$a(B)$ 在宽磁场区间显著偏离 $a_{\rm bg}$；$|\Delta B|$ 小表示共振窄而尖。$\Delta B$ 的符号决定 $a(B)$ 发散方向：$\Delta B>0$ 时 $B\to B_0^+$ 给 $a\to-\infty$，$B\to B_0^-$ 给 $a\to+\infty$。BEC 稳定性要求 $a<0$（吸引相互作用下 BEC 可存在但有限寿命），故实验上根据 $\Delta B$ 符号选择磁场扫描方向。窄共振（$|\Delta B|\ll1\,$G）需要高磁场稳定性和精度，宽共振更适合粗调但可能受多通道效应影响。

\emph{举例：$^{87}$Rb BEC 中的 Feshbach 共振。}$^{87}$Rb 在 $B_0\approx155\,$G 有广泛使用的 s 波 Feshbach 共振，$a_{\rm bg}\approx-99\,a_0$（$a_0$ 是 Bohr 半径），$\Delta B\approx0.21\,$G。在 $B=155.1\,$G 处 $a\approx+99\times(1-0.21/0.1)\approx-100\,a_0$；在 $B=155.21\,$G 处 $a=0$（无相互作用，理想 Bose 气体）；在 $B=154.9\,$G 处 $a\approx+99\times(1-0.21/(-0.1))\approx+300\,a_0$（强排斥）。通过磁场扫描，同一原子气体可从弱排斥（$a<0$，BEC 稳定）连续调谐到强排斥（$a>0$，Tonks-Girardeau 极限）或零相互作用，这是冷原子量子模拟的核心工具。

\emph{更一般结论：多通道散射理论。}实际 Feshbach 共振可能涉及多个闭通道态，此时有效 Hamiltonian 是矩阵，$a(B)$ 的极点结构由 $\det[E-H_{\rm eff}(E,B)]=0$ 的根决定。多共振重叠时，简单的 \eqnrefs{eq:feshbach-scattering-length} 失效，需要用 R-matrix 理论或耦合通道数值计算。此外，有限温度和有限密度下，平均场相互作用 $g=4\pi\hbar^2 a/m$ 使 $a(B)$ 的发散被截断，共振中心附近的有效 $a$ 饱和于 $a_{\rm eff}\sim\hbar/(m v_{\rm th})$（$v_{\rm th}$ 是热速度），这是 unitarity limit 的物理根源。

\emph{共振极点与 S 矩阵。}散射矩阵 $S(k)=\ee^{-2\ii ka}(1-2\ii k f)$ 在 $a(B)$ 发散处有简单极点。将 $a(B)=a_{\rm bg}-C_a/(B-B_0)$ 代入并解析延拓到复 $k$（$k=i\kappa$），束缚态极点 $\kappa=1/a$ 给出闭通道束缚态能量 $E_b=-\hbar^2\kappa^2/(2\mu)$。在共振附近 $E_b\propto B-B_0$，即闭通道束缚态能量与磁场线性关联，这是磁调谐 Feshbach 共振的 S 矩阵极点判据。

\emph{举例：$^{40}$K 的宽共振对比。}$^{40}$K 在 $B_0\approx202\,$G 有宽 s 波共振，$\Delta B\approx-78\,$G，$a_{\rm bg}\approx174\,a_0$。宽共振（$|\Delta B|\gg1\,$G）意味着 $a(B)$ 在大磁场范围内可调，适合研究强相互作用 Fermi 气体；但多通道效应使简单公式在 $|B-B_0|>10\,$G 时偏差增大。对比 $^{87}$Rb 的窄共振（$\Delta B\approx0.21\,$G），两种共振分别适用于不同实验目标：宽共振适合普适物理（unitarity regime），窄共振适合精细调谐相互作用零点。
\end{derivation}

\section{Landau--Zener：局部线性交叉与 Weber 方程}

现在转向含时问题。考虑 diabatic 二能级 Hamiltonian

\begin{equation}
\boxed{
H(t)
=\begin{pmatrix}
s_{\rm LZ} t/2&V\\
V^*&-s_{\rm LZ} t/2
\end{pmatrix},
}
\label{eq:lz-hamiltonian}
\end{equation}

其中 \(s_{\rm LZ}\) 的量纲是能量/时间，\(2|V|\) 是 avoided crossing 的最小绝热能隙。瞬时绝热能量

\[
E_\pm(t)
=\pm\sqrt{(s_{\rm LZ} t/2)^2+|V|^2}.
\]

\figref{fig:landau-zener} 展示了 diabatic crossing 被恒定耦合打开的过程。

\begin{figure}[tbp]
\centering
\includegraphics[width=.83\textwidth]{landau_zener.pdf}
\caption{Landau--Zener 模型的 diabatic 与 adiabatic 能级。常数耦合 \(V\) 把线性交叉打开为最小间隙 \(2|V|\)。}
\label{fig:landau-zener}
\end{figure}

通过常数相位变换可令 \(V>0\) 为实数。两分量方程为

\begin{align*}
\ii\hbar\dot c_1
&=\frac{s_{\rm LZ} t}{2}c_1+Vc_2,
\\
\ii\hbar\dot c_2
&=Vc_1-\frac{s_{\rm LZ} t}{2}c_2.
\end{align*}

消去 \(c_2\)。从第一式

\[
c_2
=\frac{\ii\hbar\dot c_1-(s_{\rm LZ} t/2)c_1}{V}.
\]

对第一式再求导并代入第二式，\(\dot c_1\) 项恰好相消，得到

\begin{equation}
\boxed{
\ddot c_1
+\left[
\frac{s_{\rm LZ}^2t^2}{4\hbar^2}
+\frac{V^2}{\hbar^2}
+\frac{\ii\,s_{\rm LZ}}{2\hbar}
\right]c_1=0.
}
\label{eq:lz-second-order}
\end{equation}

对 \(s_{\rm LZ}>0\) 定义

\[
z=\ee^{-\ii\pi/4}
\sqrt{\frac{s_{\rm LZ}}{\hbar}}t,
\qquad
\delta\equiv\frac{V^2}{\hbar s_{\rm LZ}}.
\]

由于

\[
\frac{\dd^2}{\dd t^2}
=-\ii\frac{s_{\rm LZ}}{\hbar}
\frac{\dd^2}{\dd z^2},
\qquad
 t^2=\ii\frac{\hbar}{s_{\rm LZ}}z^2,
\]

\eqnrefs{eq:lz-second-order} 变成

\begin{equation}
\boxed{
\frac{\dd^2c_1}{\dd z^2}
+\left(
\nu+\frac12-\frac{z^2}{4}
\right)c_1=0,
\qquad
\nu=-1+\ii\delta.
}
\label{eq:lz-weber}
\end{equation}

这就是 parabolic-cylinder (Weber) equation，其独立解可写为 \(D_\nu(z)\) 与适当旋转后的 \(D_\nu(-z)\)。

\section{Stokes 连接、跃迁指数与 Stückelberg 相位}

Landau--Zener 概率的非平凡之处不在“解是特殊函数”，而在如何把 \(t\to-\infty\) 的入射边界条件穿过复平面的 Stokes sectors 连接到 \(t\to+\infty\)。利用 parabolic-cylinder 函数的标准连接公式，保持 diabatic character 的概率为

\begin{equation}
\boxed{
P_{\rm LZ}
=\ee^{-2\pi\delta}
=\exp\left[-\frac{2\pi|V|^2}{\hbar|s_{\rm LZ}|}\right].
}
\label{eq:landau-zener}
\end{equation}


\begin{derivation}[Weber 函数的 Stokes 连接如何产生 Landau--Zener 指数]
\emph{第一步：Weber 方程的扇区结构。}\eqnrefs{eq:lz-weber} 的解 $D_\nu(z)$（$\nu=-1+\ii\delta$）在不同复 $z$ 扇区有不同的渐近行为。Stokes 现象是指：同一解在不同扇区的渐近展开系数不同，连接公式描述跨扇区时这些系数如何跳变。对 Landau--Zener 问题，入射边界条件在 $t\to-\infty$（$z\to\infty\cdot\ee^{-\ii\pi/4}$）给定，出射在 $t\to+\infty$（$z\to\infty\cdot\ee^{+\ii\pi/4}$），两者位于不同 Stokes 扇区，必须用连接公式把它们联系起来。

\emph{第二步：连接矩阵的振幅。}把入射/出射解按单位概率流归一化后，连接矩阵的对角振幅模为
\[
 |S_{11}|=\ee^{-\pi\delta},
\]
$\ee^{-\pi\delta}$ 的指数结构来自 parabolic-cylinder 函数 $D_\nu$ 在跨扇区时的 subdominant 指数项系数，是 Stokes 现象的精确体现。

\emph{第三步：非绝热振幅与 Gamma 函数。}非绝热转换振幅中出现的 Stokes multiplier 可写成
\[
 C_\delta
 =\frac{\sqrt{2\pi}\,\ee^{-\pi\delta/2}}
 {\sqrt{\delta}\,\Gamma(-\ii\delta)}
 \times \ee^{\ii\varphi_\delta},
\]
其中 \(\varphi_\delta\) 只影响相位。Gamma 函数 $|\Gamma(-\ii\delta)|^2$ 的模由两个恒等式决定：递推 $\Gamma(1-\ii\delta)=(-\ii\delta)\Gamma(-\ii\delta)$ 和反射公式给出的
\[
 \Gamma(1-\ii\delta)=(-\ii\delta)\Gamma(-\ii\delta),
 \qquad
 |\Gamma(1-\ii\delta)|^2
 =\frac{\pi\delta}{\sinh(\pi\delta)},
\]
后者来自 $|\Gamma(1/2+\ii y)|^2=\pi/\cosh(\pi y)$ 经函数方程推广。

\emph{第四步：概率守恒。}把 Gamma 模代入 $|C_\delta|^2$，
\begin{align*}
 |C_\delta|^2
 &=\frac{2\pi\ee^{-\pi\delta}}
 {\delta|\Gamma(-\ii\delta)|^2}\\
 &=2\ee^{-\pi\delta}\sinh(\pi\delta)\\
 &=1-\ee^{-2\pi\delta}.
\end{align*}
最后一步用 $\sinh(\pi\delta)=(\ee^{\pi\delta}-\ee^{-\pi\delta})/2$。因此单位流守恒给出
\[
 |S_{11}|^2=\ee^{-2\pi\delta},
 \qquad
 |S_{21}|^2=1-\ee^{-2\pi\delta}.
\]
前者就是保持 diabatic character 的概率\eqnrefs{eq:landau-zener}。$|S_{11}|^2+|S_{21}|^2=1$ 自动满足，验证了 Weber 方程的幺正性。指数 $\ee^{-2\pi\delta}$ 的结构来自 Gamma 函数在纯虚参量处的模，是解析数论在量子动力学中的典型体现。

\emph{Stokes 现象的物理图像。}Weber 方程的解在复 $z$ 平面上有两条 Stokes 线，沿这些线解的渐近行为发生跳变。物理上，Stokes 线对应于经典禁阻区的边界（复时间鞍点），跨线时 subdominant 指数项"突然出现"或"消失"。这一现象在 WKB 近似中也出现：连接公式正是 Stokes 现象在一维 Schrödinger 方程中的体现。Stokes multiplier 的精确值由 monodromy 决定，是复微分方程全局解结构的局部反映。

\emph{第五步：Stokes 相位。}连接系数的相位部分不能在多次 crossing 时丢弃；整理 Gamma 函数相位后正得到后文的 Stokes phase \(\phi_S\)（\eqnrefs{eq:lz-stokes-phase}）。$\phi_S$ 包含 $\arg\Gamma(1-\ii\delta)$ 的非平凡 $\delta$ 依赖，在 Stückelberg 干涉中与动力学相位叠加，产生可观测的干涉条纹。

\emph{举例：超导量子比特的 Landau--Zener 扫描。}超导 flux qubit 中，外磁通 $\Phi(t)$ 线性扫描穿过 avoided crossing，$s_{\rm LZ}=2\mu\dot\Phi$（$\mu$ 是 qubit 磁矩），$V=\Delta/2$（$\Delta$ 是能隙）。扫速 $\dot\Phi=1\,$m$\Phi_0/\mu$s、$\Delta/h=1\,$GHz 给出 $\delta=\Delta^2/(4\hbar s_{\rm LZ})\approx0.3$，绝热跟随概率 $P_{\rm LZ}=\ee^{-2\pi\delta}\approx0.15$，即 85\% 概率非绝热跳转。若 $\dot\Phi$ 减小 10 倍，$\delta\approx3$，$P_{\rm LZ}\approx\ee^{-18.8}\approx6\times10^{-9}$，几乎完全绝热。这一扫描速率依赖是量子退火和绝热量子计算的核心限制。

\emph{更一般结论。}Landau--Zener 公式 \eqnrefs{eq:landau-zener} 是 Weber 方程 Stokes 连接的最简情形。多次 avoided crossing 的 Stückelberg 干涉需要完整 $S_{\rm LZ}$ 矩阵（含 Stokes 相位）；非线性 sweep（如 $s(t)\propto t^2$）给出超椭圆积分代替 Weber 方程；多能级 crossing 需要群路径有序乘积的 Stokes 矩阵。这些推广都建立在同一复时间鞍点连接的框架上，是渐近分析在量子动力学中的典型应用。Stokes multiplier 的精确值由 Weber 方程的 monodromy 表示决定，连接了局部渐近展开与全局解的代数结构。

\emph{注记。}$\ee^{-2\pi\delta}$ 中的 $\delta=\Delta^2/(4\hbar s)$ 是能隙 $\Delta$ 与扫速 $s$ 的无量纲比，指数结构在 $s\to0$ 时给出绝热极限 $P\to0$，在 $s\to\infty$ 时给出 sudden 极限 $P\to1$。
\end{derivation}

控制参数

\[
\delta=\frac{|V|^2}{\hbar|s_{\rm LZ}|}
\]

比较“能隙平方”与“扫过速率”。\(\delta\gg1\) 时 crossing 被缓慢穿过，系统绝热跟随；\(\delta\ll1\) 时扫过太快，状态来不及混合，几乎保持原 diabatic character。

如果不仅关心单次跃迁概率，还关心多次 crossing 的 Stückelberg interference，就必须保留振幅相位。单次 Landau--Zener scattering matrix 除了 \(P_{\rm LZ}\) 外还包含 Stokes phase

\begin{equation}
\boxed{
\phi_S
=\frac\pi4
+\arg\Gamma(1-\ii\delta)
+\delta(\ln\delta-1).
}
\label{eq:lz-stokes-phase}
\end{equation}

在一种常见相位 convention 下可写

\[
S_{\rm LZ}
=\begin{pmatrix}
\sqrt{P_{\rm LZ}}\ee^{-\ii\phi_S}&-\sqrt{1-P_{\rm LZ}}\\
\sqrt{1-P_{\rm LZ}}&\sqrt{P_{\rm LZ}}\ee^{+\ii\phi_S}
\end{pmatrix}.
\]

两次 avoided crossing 之间积累的动力学相位与 \(\phi_S\) 干涉，就形成 St\"uckelberg 干涉条纹。

\begin{note}[Landau--Zener 公式何时不能直接用]
\eqnrefs{eq:landau-zener} 要求：diabatic 能差在相关时间窗内可线性化、耦合 \(V\) 近似常数、演化从远离 crossing 的 \(t\to-\infty\) 延伸到 \(t\to+\infty\)，并且系统在这段时间内可视为封闭二能级系统。有限脉冲、非线性 sweep、多个近邻能级或显著耗散都会改变连接问题，应回到一般含时传播子或多能级模型。
\end{note}

\section{Feshbach 共振与 Landau--Zener 跃迁的实验数值}\label{sec:resonance-experiments}

接下来考察超冷原子磁 Feshbach 共振。

Feshbach 共振在超冷原子物理中用于精确调控散射长度 \(a(B)\)：
\[
a(B)=a_{\rm bg}\left(1-\frac{\Delta B}{B-B_0}\right)
\]
其中 \(a_{\rm bg}\) 是背景散射长度，\(B_0\) 是共振磁场，\(\Delta B\) 是共振宽度。

\begin{center}
\begin{tabular}{lcccc}
\hline
体系 & \(B_0\) (G) & \(\Delta B\) (G) & \(a_{\rm bg}\) (\AA) & 实验组 \\
\hline
\(^6\)Li \(|1/2,-1/2\rangle\) & 834.1 & 300 & \(-1405\) & MIT (2004) \\
\(^{40}\)K \(|9/2,-9/2\rangle\) & 202.1 & 7.8 & 174 & JILA (2004) \\
\(^{87}\)Rb \(|1,1\rangle+|1,0\rangle\) & 1007.4 & 0.21 & 99 & MIT (2004) \\
\(^{23}\)Na \(|1,1\rangle\) & 907.0 & 1.0 & 52 & MIT (2003) \\
\hline
\end{tabular}
\end{center}

\begin{exbox}[\(^6\)Li 的宽共振与 BEC-BCS 渡越]
\(^6\)Li 在 834\,G 的宽共振（\(\Delta B=300\,\mathrm{G}\)）使 \(a(B)\) 可在 \(-\infty\) 到 \(+\infty\) 间连续调节。在 \(B>B_0\)（\(a>0\)）侧实现分子 BEC（紧束缚 Feshbach 分子），在 \(B<B_0\)（\(a<0\)）侧实现 BCS 超流。共振点 \(B=B_0\) 处 \(a\to\pm\infty\)，系统进入 unitary 极限：散射截面达到量子极限 \(\sigma=4\pi/k^2\)，所有长度尺度被散射长度抹平，方程只依赖密度与温度之比。MIT (2006) 在此区域观测到超流涡旋，确认了 BEC-BCS 渡越。
\end{exbox}

接下来考察Landau--Zener 跃迁的实验实现。

\begin{center}
\begin{tabular}{lccc}
\hline
体系 & \(\delta_{\rm LZ}=V^2/(\hbar|\dot\Delta|)\) & 跃迁概率 & 实验 \\
\hline
超导量子比特 (Rabi) & 0.1--10 & 0--1 可调 & Nakamura (2003) \\
Rydberg 原子 Stark & 0.5--5 & 0.1--0.9 & Gallagher (1980s) \\
分子束共振 & \(\sim1\) & 0.5 附近 & 分子物理 \\
\hline
\end{tabular}
\end{center}

\begin{exbox}[超导量子比特的 Landau--Zener sweep]
transmon 量子比特在 flux 偏置扫描中穿越 avoided crossing。能隙 \(2V/h\approx100\,\mathrm{MHz}\)，sweep 速率 \(|\dot\Delta|/h\approx1\,\mathrm{GHz}/\mathrm{ns}\)，给出 \(\delta_{\rm LZ}=V^2/(\hbar|\dot\Delta|)\approx5\)。绝热跃迁概率 \(P_{\rm LZ}=e^{-2\pi\delta_{\rm LZ}}\approx3\times10^{-14}\approx0\)——系统几乎完全跟随瞬时本征态（绝热极限）。反向快扫 \(\dot\Delta/h\approx100\,\mathrm{GHz}/\mathrm{ns}\) 给出 \(\delta_{\rm LZ}\approx0.05\)，\(P_{\rm LZ}\approx0.73\)——大部分留在初态（突然极限）。
\end{exbox}

接下来考察Fano 共振的 \(q\) 参数与谱线形。

Fano 共振的散射截面
\[
\sigma(\epsilon)=\sigma_0\frac{(q+\epsilon)^2}{1+\epsilon^2},\qquad\epsilon=\frac{E-E_r}{\Gamma/2}
\]
其中 \(q\) 是 Fano 不对称参数，\(\Gamma\) 是共振宽度。

\begin{center}
\begin{tabular}{lcc}
\hline
体系 & \(q\) & \(\Gamma\) \\
\hline
He 原子 \(2s2p\,{}^1P\) 自电离 & \(-2.8\) & 0.012 eV \\
量子点 Kondo & 0.5 & \(k_BT_K\) \\
等离子体 Fano & 1--10 & 0.1--1 eV \\
\hline
\end{tabular}
\end{center}

\begin{exbox}[He 自电离共振的 Fano 谱]
He 的 \(2s2p\,{}^1P\) 自电离共振在 60.1\,eV 附近，\(q\approx-2.8\)，\(\Gamma\approx0.012\,\mathrm{eV}\)。在 \(\epsilon=-q=2.8\) 处 \(\sigma=0\)（Fano 零点），在 \(\epsilon=1/q\approx-0.36\) 处 \(\sigma\) 达到最大。不对称线形与 Lorentzian（\(q\to\infty\)）形成鲜明对比，是离散态-连续谱干涉的直接证据。
\end{exbox}

\begin{derivation}[Fano 共振的散射相移与 Wigner--Smith 时间延迟，即\eqnrefs{eq:fano-wigner-smith-delay}]
\emph{策略路线图。}首先从 Fano 态构造提取散射矩阵 \(S(\epsilon)\) 的显式形式；其次分解 \(S\) 为背景相移与共振相移的乘积并写出共振相移 \(\delta_{\rm res}(\epsilon)\)；再次由 Wigner--Smith 公式 \(\tau=2\hbar\,\dd\delta/\dd E\) 计算时间延迟；然后分析时间延迟在 Fano 零点处的负值与因果性约束；最后与 Breit--Wigner 共振对比并给出 He 自电离的数值估计。

\emph{第一步：Fano 态与 S 矩阵。}Fashbach 投影给出离散态 \(|\phi\rangle\) 与连续态 \(|E\rangle\) 的耦合。Fano 对角化产生 Fano 态
\begin{equation}
|\Psi_E\rangle=\frac{q+\epsilon}{\sqrt{\epsilon^2+1}}\,|\phi\rangle+\frac{1}{\sqrt{\epsilon^2+1}}\,|E\rangle,
\qquad
\epsilon=\frac{E-E_r}{\Gamma/2},
\label{eq:fano-state-s-matrix}
\end{equation}
即\eqnrefs{eq:fano-normalized-state} 的参数化。散射矩阵元 \(S(E)=\langle E_{\rm out}|\Psi_E\rangle/\langle E_{\rm in}|\Psi_E\rangle\) 由入射连续态到出射连续态的跃迁确定。利用 Feshbach 有效 Hamiltonian（即\eqnrefs{eq:feshbach-effective}）在单共振近似下的形式，S 矩阵分解为
\begin{equation}
S(\epsilon)=\ee^{2\ii\delta_{\rm bg}}\,\frac{\epsilon+q-\ii}{\epsilon+q+\ii},
\label{eq:fano-s-matrix}
\end{equation}
其中 \(\delta_{\rm bg}\) 是背景相移（来自势散射和远程连续态贡献），\((\epsilon+q-\ii)/(\epsilon+q+\ii)\) 是共振极点结构。注意\eqnrefs{eq:fano-s-matrix} 的极点在 \(\epsilon=-q-\ii\)（下半平面），保证因果性。

\emph{第二步：共振相移。}其次从 \(S=\ee^{2\ii\delta}\) 提取总相移 \(\delta=\delta_{\rm bg}+\delta_{\rm res}\)。共振相移
\begin{equation}
\delta_{\rm res}(\epsilon)=-\arctan\!\frac{1}{\epsilon+q},
\label{eq:fano-resonance-phase-shift}
\end{equation}
即\eqnrefs{eq:fano-phase-shift}。当 \(\epsilon\) 从 \(-\infty\) 扫到 \(+\infty\)，\(\delta_{\rm res}\) 从 \(0\) 单调增到 \(\pi\)（经过 \(\epsilon=-q\) 处 \(\delta_{\rm res}=\pi/2\)）。总相移跃变 \(\Delta\delta=\pi\) 是共振的拓扑标志——与 Breit--Wigner 共振（\(q\to\infty\)，\(\delta_{\rm res}=-\arctan(1/\epsilon)\)）相同。

\emph{第三步：Wigner--Smith 时间延迟。}再次定义 Wigner--Smith 时间延迟
\begin{equation}
\tau(\epsilon)=2\hbar\,\frac{\dd\delta}{\dd E}
=\frac{4\hbar}{\Gamma}\,\frac{\dd\delta_{\rm res}}{\dd\epsilon},
\label{eq:wigner-smith-definition}
\end{equation}
因子 \(4\hbar/\Gamma\) 来自 \(\dd E=(\Gamma/2)\dd\epsilon\)。由\eqnrefs{eq:fano-resonance-phase-shift}，
\begin{equation}
\frac{\dd\delta_{\rm res}}{\dd\epsilon}
=\frac{1}{1+(\epsilon+q)^2},
\end{equation}
故
\begin{equation}
\tau(\epsilon)=\frac{4\hbar}{\Gamma}\,\frac{1}{1+(\epsilon+q)^2}.
\label{eq:fano-wigner-smith-delay}
\end{equation}
即\eqnrefs{eq:fano-wigner-smith-delay}。时间延迟在 \(\epsilon=-q\)（Fano 零点）处取极大值 \(\tau_{\max}=4\hbar/\Gamma\)，与 Breit--Wigner 结果相同——共振中心粒子被离散态暂时俘获，延迟约 \(\hbar/\Gamma\) 量级（共振寿命）。

\emph{第四步：负时间延迟与因果性。}然后注意\eqnrefs{eq:fano-wigner-smith-delay} 对所有 \(\epsilon\) 为正——Fano 共振的 Wigner--Smith 延迟非负。但更精细的分析考虑背景相散 \(\delta_{\rm bg}\) 的能量依赖：若 \(\dd\delta_{\rm bg}/\dd E<0\)（反常色散区，如势垒散射），总时间延迟
\begin{equation*}
\tau_{\rm total}=\frac{4\hbar}{\Gamma}\,\frac{1}{1+(\epsilon+q)^2}
+2\hbar\,\frac{\dd\delta_{\rm bg}}{\dd E}
\end{equation*}
可在远离共振处为负。负时间延迟（"时间超前"）不违反因果性：Wigner--Smith 公式给出的是波包中心延迟，而非信号前沿延迟。因果性要求信号前沿不超前，由 Paley--Wiener 定理保证（S 矩阵解析性）。Fano 零点 \(\epsilon=-q\) 处 \(\sigma=0\) 但 \(\tau=\tau_{\max}\)：粒子被完全俘获后以特定相位重新发射，干涉相消使截面为零但延迟仍为正。

\emph{第五步：He 自电离的数值估计。}最后对 He \(2s2p\,{}^1P\) 自电离共振，\(\Gamma\approx0.012\,\si{eV}\)，\(q\approx-2.8\)。最大时间延迟
\begin{equation*}
\tau_{\max}=\frac{4\hbar}{\Gamma}
=\frac{4\times6.58\times10^{-16}\,\si{eV\cdot s}}{0.012\,\si{eV}}
\approx2.2\times10^{-13}\,\si{s}=220\,\si{fs}.
\end{equation*}
在 Fano 零点 \(\epsilon=2.8\)（\(E\approx60.1+0.017\,\si{eV}\)）处，He 原子被自电离态俘获约 220\,fs 后重新发射电子。对比直接光电离时间（\(\sim1\,\si{fs}\)），共振延迟使自电离过程慢约 200 倍——这是 attosecond 科学中 Fano 共振时间分辨实验的基础。


\emph{[Fano 共振的解析结构与 Fano--Feshbach 耦合]} Fano 共振的线型由离散态与连续态的耦合给出，其解析结构包含共振位置、宽度和 Fano 参数 \(q\)。分四步建立。

\emph{第一步：Fano 模型。}Hamiltonian \(H=H_0+V\)，\(H_0|\phi\rangle=E_0|\phi\rangle\)（离散态），\(H_0|E\rangle=E|E\rangle\)（连续态），\(V\) 耦合离散与连续：\(V_\phi(E)=\langle\phi|V|E\rangle\)。Fano 假设 \(V_\phi(E)\) 在共振附近缓变（常数近似），自能 \(\Sigma(E)=\Delta(E)-i\Gamma(E)/2\)，\(\Gamma=2\pi|V_\phi(E_0)|^2\)（衰变率），\(\Delta=\mathcal P\int\frac{|V_\phi(E')|^2}{E_0-E'}\dd E'\)（能移）。

\emph{第二步：Fano 线型。}散射振幅 \(f(E)=\frac{q+\epsilon}{\epsilon+i}\)（\(\epsilon=(E-E_r)/(\Gamma/2)\)，\(E_r=E_0+\Delta\) 共振能，\(q\) 是 Fano 参数）。截面
\begin{equation}
  \sigma(E)=\sigma_0\frac{(q+\epsilon)^2}{1+\epsilon^2},
  \label{eq:phys-rl-fano-profile}
\end{equation}
其中 \(\sigma_0\) 是连续态背景截面。\(q\to\infty\) 给出 Lorentzian（纯离散共振）；\(q=0\) 给出反共振（窗口共振，\(\sigma\propto\epsilon^2/(1+\epsilon^2)\)，共振处完全抑制）；一般 \(q\) 给出不对称线型。Fano 参数 \(q=-\frac{\mathcal P\int V_\phi(E')/(E_0-E')\dd E'}{\pi V_\phi(E_0)}\) 是离散态在连续态中的相位移——纯几何量。

\emph{第三步：Feshbach 共振与磁场调控。}在冷原子物理中，Feshbach 共振通过磁场调控散射长度：\(a(B)=a_{\rm bg}\left(1-\frac{\Delta B}{B-B_0}\right)\)（\(a_{\rm bg}\) 背景散射长度，\(B_0\) 共振磁场，\(\Delta B\) 共振宽度）。物理机制：开通道（碰撞态）与闭通道（分子束缚态）通过超精细耦合混合，磁场改变 Zeeman 能移使闭通道与开通道阈值交叉。在 \(B=B_0\) 处 \(a\to\infty\)（unitary 极限），分子结合能在 \(B\) 扫过 \(B_0\) 时从 \(E<0\)（分子）连续变到 \(E>0\)（虚态）——Feshbach 共振是冷原子物理中控制相互作用强度的旋钮。

\emph{第四步：多通道 Fano--Feshbach 理论。}多通道散射中 Fano 共振的推广：开通道 \(\{|\chi_\alpha\rangle\}\) 和闭通道 \(\{|\phi_n\rangle\}\) 的耦合矩阵 \(W_{\alpha n}=\langle\chi_\alpha|V|\phi_n\rangle\)。多通道 Fano 参数 \(\boldsymbol q\) 是向量，截面是各开通道贡献的 Fano 线型叠加。共振位置由闭通道本征值（含自能修正）的零点给出，宽度由耦合矩阵的奇异值给出。多通道理论统一处理冷原子 Feshbach 共振（核-电子超精细耦合）、原子物理自电离共振（电子-连续态耦合）和核物理共振（复合核-散射态耦合），是共振物理的普适框架。
\end{derivation}
